% ============================================================
% CUMCM 国赛论文 LaTeX 模板（完整骨架 + 详细注释）
% 编译：XeLaTeX（cumcmthesis.cls 强制要求 xelatex）
% 用法：复制本文件 + templates/cumcmthesis.cls 到 paper/ 目录
%      图片放 ../figures/，数据放 ../data/
% ============================================================

\documentclass[withoutpreface,bwprint]{cumcmthesis}
% withoutpreface：不生成承诺书页（正式提交由官方系统生成）
% bwprint：黑白打印（避免高饱和彩色被扣分）

% ---- 宏包（cumcmthesis.cls 已内置 booktabs/listings/subcaption 等）----
\usepackage{url, subcaption, geometry, setspace, array, graphicx,
            multirow, hyperref, amsmath, amssymb, bm, longtable,
            booktabs, float}
% booktabs=三线表 | float=[H]强制位置 | amsmath=公式

\graphicspath{{../figures/}}   % 图片统一指向 figures/ 目录
\geometry{a4paper,top=2.54cm,bottom=2.54cm,left=3.18cm,right=3.18cm}  % 官方规范：上下2.54cm、左右3.18cm

\title{论文标题——副标题点明核心方法}

\begin{document}
\maketitle

% ============================================================
% 摘要（单独成页，≤1页且尽量占满，最后写）
% 铁律：每问必须给【具体数值】+【核心公式】+【结果讨论】
%       禁止"结果较好/效果显著"，禁止"针对问题n"套话开头
% ============================================================
\begin{abstract}
第1段（2-3句总起）：本文针对[问题]，建立了[模型体系]，采用[方法框架]，
并通过[验证方式]论证了可靠性。

第2-N段（逐问一段，每问 = 数学抽象→方法→关键数值→验证）：
\textbf{问题一（模型名）：}基于[原理]，推导了[核心公式]，得到[具体数值]。

\textbf{问题二（模型名）：}设计[算法]，得到
\textbf{$7.58\pm0.62\,\mu\text{m}$}（相对偏差$1.7\%$，$R^2=0.9996$）。

末段（2-3句）：模型评价 + 创新点 + 推广价值。

\keywords{\textbf{[关键词1]\quad [关键词2]\quad [关键词3]\quad [关键词4]\quad [关键词5]}}
% 关键词 3-5 个，按重要性排序，禁止引号、禁止编号
\end{abstract}

\tableofcontents   % 目录（可选）

% ============================================================
% 一、问题重述
% ============================================================
\section{问题重述}
\subsection{问题背景}
% 工程/实际价值 + 简要原理，1-2 段
\subsection{问题提出}
% 用 itemize 列出题目各问，编号 1/2/3，不要照抄原文要提炼

% ============================================================
% 二、问题分析
% ============================================================
\section{问题分析}
% 逐问分析思路：每问的数学本质归类 + 方法选择理由
% 末尾放"技术路线总图"（fig_pipeline）收束

% ============================================================
% 三、模型假设与符号说明
% ============================================================
\section{模型假设与符号说明}
\subsection{模型假设}
% enumerate 每条：假设内容 + 成立理由 + 影响评估
% 禁止"忽略其他次要因素"这类模糊表述，要具体
\subsection{符号说明}
% 三线表分组列出（见下方三线表示例）

% ============================================================
% 四、模型的建立与求解（核心，逐问展开）
% ============================================================
\section{模型的建立与求解}

\subsection{问题一：模型名称}
\subsubsection{模型建立}
% 公式用 equation 环境（自动编号），紧跟"物理意义解读"段落
% \begin{equation}\label{eq:xxx}
%   ... 公式 ...
% \end{equation}
\subsubsection{算法设计}
% enumerate 列算法步骤，说明选型理由 + 替代方案被否理由
\subsubsection{结果分析}
% 三线表汇总 + 结果解读（解释方法间差异来源）
\subsubsection{模型检验}
% 交叉验证/误差分析，引用图表

\subsection{问题二：模型名称}
% 同上结构

\subsection{问题三：模型名称}
% 同上结构

% ============================================================
% 五、灵敏度分析（必做）
% ============================================================
\section{灵敏度分析}
% 改变关键参数 ±X%，看结果稳定性；配龙卷风图/灵敏度表

% ============================================================
% 六、模型评价与改进
% ============================================================
\section{模型评价与改进}
\subsection{模型优点}
% enumerate
\subsection{模型不足与改进方向}

% ============================================================
% 七、结论
% ============================================================
\section{结论}
% 分条 (1)(2)(3)，每条对应一问题 + 具体数值

% ============================================================
% 三线表示例（booktabs，全项目统一，禁止全边框）
% ============================================================
% \begin{table}[H]
% \centering
% \caption{表题（表题在表上方）}
% \label{tab:result}
% \begin{tabular}{ccccc}
% \toprule
% 方法 & 附件1 & 附件2 & 平均值 & 相对偏差 \\
% \midrule
% 方法A & $7.53\,\mu\text{m}$ & $7.52\,\mu\text{m}$ & $7.53\,\mu\text{m}$ & $0.1\%$ \\
% 方法B & $7.39\,\mu\text{m}$ & $7.63\,\mu\text{m}$ & $7.51\,\mu\text{m}$ & $3.2\%$ \\
% \midrule
% \textbf{综合} & \textbf{$7.51\,\mu\text{m}$} & \textbf{$7.64\,\mu\text{m}$} & \textbf{$7.58\,\mu\text{m}$} & \textbf{$1.7\%$} \\
% \bottomrule
% \end{tabular}
% \end{table}

% ============================================================
% 图插入示例（[H]强制位置，图题在图下方，caption 标 (a)(b)(c)(d)）
% ============================================================
% \begin{figure}[H]
% \centering
% \includegraphics[width=0.95\textwidth]{fig7_data_processing.png}
% \caption{图题直接点结论：(a)原始光谱与SG平滑对比；(b)...}
% \label{fig:data_processing}
% \end{figure}

% ============================================================
% 参考文献（不编号为"第X部分"，thebibliography 手写）
% GB/T 7714 著录格式：
%   期刊 [J]：作者. 题名[J]. 刊名, 年, 卷(期): 起止页码.
%   专著 [M]：作者. 书名[M]. 版本. 出版地: 出版者, 年: 页码.
%   标准 [S]：标准编号, 标准名称[S]. 出版地: 出版者, 年.
%   AI工具：工具名, 版本, 开发机构, 使用日期.
% ============================================================
% \begin{thebibliography}{20}
% \bibitem{ref1} 作者A, 作者B. 论文题名[J]. 期刊名, 2024, 12(3): 45-52.
% \bibitem{ref2} 作者C. 书名[M]. 2版. 北京: 出版社, 2023: 100-120.
% \bibitem{ref_ai} Claude, Opus 4, Anthropic, 2026年7月.
% \end{thebibliography}

% ============================================================
% 附录（核心算法伪代码 + 程序清单，不放真实代码）
% ============================================================
% \section*{附录}
% \subsection*{附录A：核心算法伪代码}
% % verbatim 环境放伪代码（不编号）
% \subsection*{附录B：程序文件清单}
% % 三线表列"文件名 + 功能 + 运行依赖"

\end{document}
